\documentclass[letterpaper,10pt,conference]{ieeeconf}
\IEEEoverridecommandlockouts
\overrideIEEEmargins
\usepackage{xcolor}

\usepackage{amsmath,amssymb,amsthm,mathtools,bm}
\usepackage[hidelinks]{hyperref}
\usepackage{enumitem}
\usepackage{microtype}
\usepackage{cite}
\allowdisplaybreaks
\emergencystretch=1.2em
\usepackage{graphicx}
\newtheorem{assumption}{Assumption}
\newtheorem{definition}{Definition}
\newtheorem{proposition}{Proposition}
\newtheorem{theorem}{Theorem}
\newtheorem{remark}{Remark}
\usepackage{algorithm}
\usepackage{algpseudocode}
\newcommand{\R}{\mathbb{R}}
\newcommand{\N}{\mathbb{N}}
\newcommand{\PP}{\mathbb{P}}
\newcommand{\Bm}{B_{\gamma}^{\tilde f}}
\newcommand{\Bt}{B_{\gamma}^{f}}
\newcommand{\fapp}{\widetilde{f}}
\newcommand{\safe}{\mathcal{L}}
\newcommand{\fail}{\mathcal{K}}
\newcommand{\xv}{x}
\newcommand{\uv}{u}
\newcommand{\dv}{d}
\newcommand{\Uv}{\mathcal{U}}
\newcommand{\Dv}{\mathcal{D}}
\newcommand{\Xiv}{\Xi}
\newcommand{\Xrl}{\mathcal{X}}
\newcommand{\norm}[1]{\left\lVert #1 \right\rVert}
\newcommand{\ip}[2]{\left\langle #1,#2 \right\rangle}
\newcommand{\pos}[1]{\left[#1\right]_+}

\title{Conformal Control Barrier--Value Function Shielding under Approximate Dynamics}
\author{}

\begin{document}
\maketitle

\begin{abstract}
% To be added.
\end{abstract}

\section{Introduction}
%check abbbrivation being used just once and p
%safe RL (cite works in rl and explain)
%without dynamics (model and environment) it may become unsafe , then works that have been done on how to make rl safe, how safty filters work and cite, how safty filters work with known dynamic and cie, also cite cbf, or hj, 
%explanations on cbf and hj then limitations and cbvf (but just ref to some works and not too deep since i talk about them later ) (connect about needing interaction with theword and dynamic may not be correct) and limitations that led to my work
%what this work is about and cbvf and novelty and compare to micheles
%what each section will be 


%\subsection{Limitations of Interaction-Only Safety Learning}

%A common line of work connects reachability-based safety analysis with
%reinforcement learning and suggests that safety certificates can, in principle,
%be learned from interaction data \cite{fisac2019general,lu2025adaptiveshielding}.
%This perspective is appealing because it reduces reliance on a complete
%analytical model of the system and instead attempts to recover a safety value
%from experience.

%In safety-critical settings, however, relying only on interaction can be
%difficult. Learning an accurate reachability-style safety object may require
%large amounts of data, and early exploration can risk trajectories that pass
%near the boundary of the failure set $\fail$ defined in
%\eqref{eq:safe_failure_sets}. Avoiding such events often requires conservative
%initializations or restrictive exploration, which can limit the useful data
%available for learning.

%For the setting considered here, it is therefore more practical to begin with a
%model-based certificate and then correct its safety-relevant error. This shifts
%the burden from learning the entire safety object online to calibrating only the
%part of the mismatch that can weaken the shielding condition. Intuitively,
%model errors that move the system away from the unsafe set do not require
%additional conservatism; the critical errors are those that make the certificate
%decrease more than the approximate model predicts.
%%%%%%%%%%%
%Existing work connects reachability-based safety analysis with reinforcement
%learning and shows that safety certificates can, in principle, be learned from
%interaction data \cite{fisac2019general,lu2025adaptiveshielding}. However, in
%safety-critical settings, learning such certificates only from interaction can
%require substantial data and may expose the system to unsafe behavior during
%exploration. Instead of learning a new certificate entirely from data, this
%paper starts from a model-based CBVF certificate and calibrates the part of the
%model mismatch that can weaken its online safety condition.

\section{Problem Formulation and Preliminaries}

This section presents the problem formulation used in the proposed conformal
CBVF shielding framework. We first describe the safe reinforcement-learning
problem in discrete time, where a policy is trained for task performance while
accounting for a safety-related cost. Since this expected-cost formulation does
not by itself prevent unsafe states along individual trajectories, it can result
in occasional safety violations during deployment. The learned policy is
therefore used as a nominal controller inside a continuous-time shielding
framework. We then introduce the continuous-time safety model, including the true dynamics,
the approximate dynamics available to the shield, and the target function that
defines the safe and failure sets. Using this model, we formulate the control
barrier--value function and derive the associated online safety filter. This
filter provides the model-based baseline that will later be calibrated to
account for mismatch between the approximate and true dynamics.

%%%%%%%%%%%%%%%%%%%%%%
\subsection{Safe reinforcement learning and the shielding interface}

The learning problem is modeled as a discounted constrained Markov decision
process (CMDP) \cite{altman1999cmdp,sutton2018rl}. Let
$\Xrl\subseteq\R^n$ denote the state space, and let
$\Uv\subseteq\R^m$ denote the action space. A discounted CMDP is specified by
\begin{equation}
    (\Xrl,\Uv,P,r,c,\rho,\lambda,d_c),
    \label{eq:cmdp_tuple}
\end{equation}
where $P$ is the transition kernel, $r:\Xrl\times\Uv\to\R$ is the reward
function, $c:\Xrl\times\Uv\to\R_{\ge 0}$ is a safety-related cost,
$\rho\in\mathcal{P}(\Xrl)$ is the initial-state distribution,
$\lambda\in[0,1)$ is the discount factor, and $d_c\in\R_{\ge 0}$ is the cost
budget. Let $\pi:\Xrl\to\mathcal{P}(\Uv)$ denote a policy learned by the RL
agent. The discounted return and discounted cost under $\pi$ are
\begin{equation}
\begin{aligned}
    J_r^\pi(\rho)
    &:=
    \mathbb{E}_{\rho}^{\pi}\!\left[\sum_{k=0}^{\infty}\lambda^k r(x_k,u_k)\right],\\
    J_c^\pi(\rho)
    &:=
    \mathbb{E}_{\rho}^{\pi}\!\left[\sum_{k=0}^{\infty}\lambda^k c(x_k,u_k)\right].
\end{aligned}
\label{eq:cmdp_objectives}
\end{equation}
The CMDP objective is to maximize $J_r^\pi(\rho)$ subject to
\begin{equation}
    J_c^\pi(\rho)\le d_c.
    \label{eq:cmdp_constraint}
\end{equation}
In practice, CMDP objectives of the form
\eqref{eq:cmdp_objectives}--\eqref{eq:cmdp_constraint} are commonly optimized
using constrained policy-search or Lagrangian safe-RL methods
\cite{achiam2017cpo,tessler2019rcpo}. The shielding framework below does not
depend on a particular safe-RL optimizer; it only uses the resulting policy as
the nominal controller.

The CMDP formulation captures safety through the expected cumulative cost in
\eqref{eq:cmdp_objectives} and the budget constraint in
\eqref{eq:cmdp_constraint}. However, this constraint limits the discounted cost
in expectation, rather than enforcing safety at every state along each realized
trajectory. Thus, a policy that performs well in the CMDP sense can still
produce unsafe rollouts during deployment.

This limitation motivates a distinction between the discrete-time learning
problem and the continuous-time safety model used below. The CMDP formulation
in \eqref{eq:cmdp_tuple}--\eqref{eq:cmdp_constraint} describes the nominal
learning problem at decision indices $k$, while the safety analysis is carried
out for continuous-time trajectories. Along a rollout, the sampled CMDP state
is identified with the continuous-time state at the corresponding decision
time, i.e., $x_k=\xv(t_k)$. The action selected by the RL policy is held
constant between decision times, unless it is modified by the online safety
filter. This zero-order-hold connection provides the interface between the
discrete-time learner and the continuous-time system.

At deployment, the stochastic policy $\pi$ is evaluated deterministically, for
example by using its mean action. We denote the resulting nominal controller by
\begin{equation}
    \pi_{\mathrm{nom}}:\R^n\times[t_0,0]\to\Uv,
    \qquad t_0\in\R_{<0},
    \label{eq:nominal_policy}
\end{equation}
where the interval $[t_0,0]$ is the finite backward-time horizon used in the
reachability-based safety analysis. The action applied to the system is then
given by
\begin{equation}
    \uv(s)=\phi\!\left(\xv(s),s,\pi_{\mathrm{nom}}(\xv(s),s)\right),
    \qquad s\in[t_0,0],
    \label{eq:shield_map}
\end{equation}
where $\phi:\R^n\times[t_0,0]\times\Uv\to\Uv$ is the online safety-filter map.

Thus, the RL policy provides the nominal action for task performance, while
the online safety filter modifies this action only when needed for safety. The
next subsection makes this interface precise by introducing the continuous-time
dynamics and safety objective used by the filter.




%%%%%%%%%%%%%%%%%%%% model 
\subsection{Continuous-time safety model}

We now introduce the continuous-time model used for safety analysis. The true
system evolves over a finite horizon $[t_0,0]$, with $t_0<0$, according to
\begin{equation}
    \dot{\xv}(s)=f\bigl(\xv(s),\uv(s),\dv(s)\bigr),
    \qquad s\in[t_0,0],
    \label{eq:true_dynamics}
\end{equation}
where
\begin{equation}
    \xv(s)\in\R^n,
    \qquad
    \uv(s)\in\Uv,
    \qquad
    \dv(s)\in\Dv,
    \qquad
    \Dv\subset\R^p.
    \label{eq:spaces}
\end{equation}
Here $\Dv$ is a compact disturbance set, and the input set $\Uv$ is as
introduced in the CMDP formulation in \eqref{eq:cmdp_tuple}.

The safety filter is constructed from an approximate model rather than from the
true dynamics in \eqref{eq:true_dynamics}. We therefore also consider
\begin{equation}
    \dot{\xv}(s)=\fapp\bigl(\xv(s),\uv(s),\dv(s)\bigr).
    \label{eq:model_dynamics}
\end{equation}
The same disturbance set $\Dv$ is used in \eqref{eq:true_dynamics} and
\eqref{eq:model_dynamics}. The discrepancy between $f$ and $\fapp$ represents
the safety-relevant modeling error addressed by the calibration step in
Section~\ref{sec:conformal_calibration}. This modeling error can arise from parametric uncertainty, unmodeled dynamics,
or other discrepancies between the approximate model and the true system.

We use $\Uv[t_0,0]$ and $\Dv[t_0,0]$ to denote the sets of Lebesgue-measurable
control and disturbance signals on $[t_0,0]$. As in robust reachability, the
disturbance may act through a nonanticipative strategy
$\xi_{\dv}:\Uv[t_0,0]\to\Dv[t_0,0]$, and $\Xiv[t_0,0]$ denotes the set of such
strategies \cite{choi2021cbvf,mitchell2005time}.

Safety is encoded by a target function $l:\R^n\to\R$, and the safe and failure
sets are defined by
\begin{equation}
\begin{aligned}
    \safe &:= \{\xv\in\R^n : l(\xv)\ge 0\},\\
    \fail &:= \{\xv\in\R^n : l(\xv)<0\}.
\end{aligned}
\label{eq:safe_failure_sets}
\end{equation}
The goal is to keep the true trajectory of \eqref{eq:true_dynamics} inside
$\safe$ throughout the horizon.

\begin{assumption}[Standing regularity]
\label{ass:regularity}
The functions $f$ and $\fapp$ are continuous in $(\xv,\uv,\dv)$, locally
Lipschitz in $\xv$ uniformly over $(\uv,\dv)\in\Uv\times\Dv$, and bounded on
compact sets. The target function $l$ is bounded, locally Lipschitz, and
continuously differentiable on the region where terminal gradients are
evaluated.
\end{assumption}

Assumption~\ref{ass:regularity} guarantees existence of trajectories on finite
horizons and provides the regularity needed for the reachability-based
certificate developed below. With the continuous-time safety model in place, we
next formulate the model-based certificate used by the online safety filter.




%%%%% cbvf 
\subsection{Model-based control barrier--value function}

The safety filter is built from a CBVF
associated with the approximate model \eqref{eq:model_dynamics}. The CBVF
construction uses the finite horizon $[t_0,0]$, where $0$ denotes the terminal
time of the reachability lookahead. The reachability computation is posed
backward from this terminal time, while trajectories evolve forward in the time
variable from $t_0$ to $0$.

CBVFs connect Hamilton--Jacobi reachability with barrier-style online filtering
\cite{choi2021cbvf}. For a discount parameter $\gamma\in\R_{\ge 0}$, define
\begin{equation}
\begin{aligned}
    \Bm(\xv,t)
    :=
    \inf_{\xi_{\dv}\in\Xiv[t,0]}
    \sup_{\uv(\cdot)\in\Uv[t,0]}
    \min_{r\in[t,0]}
    e^{\gamma(r-t)}l\bigl(\tilde{\xv}(r)\bigr),
\end{aligned}
\label{eq:model_cbvf}
\end{equation}
where $t\in[t_0,0]$. When $\gamma=0$, \eqref{eq:model_cbvf} reduces to the
standard finite-horizon reachability value function. For $\gamma>0$, it gives
the discounted barrier-value formulation of CBVFs \cite{choi2021cbvf}.

In \eqref{eq:model_cbvf}, the state satisfies $\tilde{\xv}(t)=\xv$, and
$\tilde{\xv}(\cdot)$ evolves according to
\begin{equation}
    \dot{\tilde{\xv}}(r)
    =
    \fapp\bigl(
    \tilde{\xv}(r),
    \uv(r),
    \xi_{\dv}[\uv](r)
    \bigr),
    \qquad r\in[t,0].
    \label{eq:model_cbvf_traj}
\end{equation}
The trajectory equation \eqref{eq:model_cbvf_traj} uses the approximate
dynamics in \eqref{eq:model_dynamics} and the nonanticipative disturbance
strategy introduced above. Under Assumption~\ref{ass:regularity} and the
comparison results for Hamilton--Jacobi reachability variational inequalities,
$\Bm$ is characterized as the unique viscosity solution of
\begin{equation}
\label{eq:model_cbvf_vi}
\begin{aligned}
0=\min\Bigl\{\, &l(\xv)-\Bm(\xv,t),\\
&\partial_t \Bm(\xv,t) + \gamma \Bm(\xv,t) \\
&\quad + \max_{\uv\in\Uv}\min_{\dv\in\Dv}
\ip{\nabla_{\xv}\Bm(\xv,t)}{\fapp(\xv,\uv,\dv)}
\Bigr\},
\end{aligned}
\end{equation}
with terminal condition
\begin{equation}
    \Bm(\xv,0)=l(\xv).
    \label{eq:model_cbvf_terminal}
\end{equation}
This viscosity characterization follows from the standard Hamilton--Jacobi
reachability construction together with the CBVF formulation
\cite{mitchell2005time,choi2021cbvf}.

Equation \eqref{eq:model_cbvf_vi} yields a property used throughout the paper:
\begin{equation}
    \Bm(\xv,t)\le l(\xv),
    \qquad \forall (\xv,t)\in\R^n\times[t_0,0],
    \label{eq:Bm_le_l}
\end{equation}
which implies that any point satisfying $\Bm(\xv,t)\ge 0$ is safe with respect
to the target function $l$ in \eqref{eq:safe_failure_sets}. Accordingly, define
the zero-superlevel set
\begin{equation}
    \mathcal{C}_{\gamma}^{\tilde f}(t)
    :=
    \{\xv\in\R^n : \Bm(\xv,t)\ge 0\},
    \label{eq:model_safe_set}
\end{equation}
which is the certified safe region induced by the approximate model
\eqref{eq:model_dynamics}. This set provides the model-based certificate
enforced by the safety filter. The next subsection shows how this certificate
leads to an online filtering problem.


%%%%%%%%%%%%%%% safety filter
\subsection{Model-based safety filter}

The model-based certificate in \eqref{eq:model_safe_set} can be enforced online
by filtering the nominal action in \eqref{eq:nominal_policy}. At each state and
time, the filter selects the closest admissible input to the nominal action
while satisfying the CBVF decrease condition induced by the approximate model
\eqref{eq:model_dynamics}:
\begin{equation}
\label{eq:model_cbvf_qp}
\begin{aligned}
\pi_{\mathrm{CBVF}}^{\tilde f}(\xv,t)
=\;& \arg\min_{\uv\in\Uv}
\norm{\uv-\pi_{\mathrm{nom}}(\xv,t)}^2 \\
\text{s.t.}\;&
\partial_t \Bm(\xv,t) + \gamma \Bm(\xv,t) \\
&\quad + \min_{\dv\in\Dv}
\ip{\nabla_{\xv}\Bm(\xv,t)}{\fapp(\xv,\uv,\dv)}
\ge 0.
\end{aligned}
\end{equation}
Under bounded input constraints, feasibility of \eqref{eq:model_cbvf_qp} can
depend on the current state and time. The filter is therefore interpreted on
the region where the constraint admits at least one admissible input.

When the approximate model \eqref{eq:model_dynamics} is control-affine in
$\uv$ and $\Uv$ is a compact convex input set, such as a box or polytope of
actuator limits, the constraint in \eqref{eq:model_cbvf_qp} is affine in
$\uv$. The online problem is then a quadratic program. As in standard
CBF-based safety filters, the optimizer returns the nominal input whenever
$\pi_{\mathrm{nom}}(\xv,t)$ already satisfies the safety condition
\cite{ames2019cbf,choi2021cbvf}.

The filter is implemented at decision times
\begin{equation}
    t_0<t_1<\cdots<t_H=0,
    \label{eq:decision_times}
\end{equation}
where $H\in\N$, and the chosen control is held constant on each interval
$[t_j,t_{j+1})$. This provides the interface to the nominal controller in
\eqref{eq:nominal_policy}. Since the derivative condition in
\eqref{eq:model_cbvf_qp} is enforced only at decision times, the zero-order-hold
implementation introduces a sampled-data gap between the continuous-time
analysis and the implemented filter.

This formulation remains model-dependent, since the certificate enforced in
\eqref{eq:model_cbvf_qp} is constructed from the approximate dynamics
\eqref{eq:model_dynamics} rather than the true dynamics
\eqref{eq:true_dynamics}. The next section addresses this gap by calibrating
the safety-relevant mismatch between the two models.




%%%%%%%%%%%%%%%%%
\section{Safety Certification under Model Mismatch}

This section explains how model mismatch affects the model-based CBVF filter.
The baseline is the CBVF certificate constructed from the approximate model,
while safety is required for the true closed-loop system. Rather than learning a
new safety certificate entirely from interaction data, we keep the model-based
CBVF structure and calibrate the part of the mismatch that can make the online
safety condition too optimistic.

The main idea is to use a data-driven safety buffer inside the CBVF filter. This
buffer tightens the online constraint by an amount calibrated from rollout data.
The resulting calibrated filter will later yield a probabilistic safety
guarantee for the true trajectory over the finite horizon.

\subsection{Safety-relevant model mismatch}

The disturbance set $\Dv$ in \eqref{eq:spaces} captures exogenous variability
that is explicitly represented in the true dynamics \eqref{eq:true_dynamics}
and the approximate dynamics \eqref{eq:model_dynamics}. This should be
distinguished from modeling error between $f$ and $\fapp$. One could enlarge
$\Dv$ to absorb such error, but doing so can make the offline reachability
computation overly conservative and can substantially shrink the certified set
$\mathcal{C}_{\gamma}^{\tilde f}(t)$ in \eqref{eq:model_safe_set}. On the other
hand, if the modeling error is not accounted for, the certified set induced by
the approximate model need not be forward invariant for the true system.

The relevant error for the online filter is not the full discrepancy between
$f$ and $\fapp$, but the component that enters the CBVF constraint in
\eqref{eq:model_cbvf_qp}. In particular, the safety condition becomes
optimistic when the directional derivative of the certificate along the true
dynamics is smaller than the derivative predicted by the approximate model.
Errors in the opposite direction do not require additional conservatism because
they make the certificate decrease less than predicted.

This motivates a one-sided correction that targets only the safety-degrading
directional mismatch. A fixed worst-case buffer can be overly conservative,
whereas no correction can leave the model-based filter unsafe under mismatch.
Before introducing the conformal buffer, we first describe a conservative
set-tightening relation that illustrates what a value-function-level correction
would require.

\subsection{A conservative set-tightening relation}

A conservative way to account for model mismatch is to tighten the model-based
certificate. A simple comparison argument shows what such a tightening would
require at the value-function level. Let $\Bt$ denote the CBVF defined as in
\eqref{eq:model_cbvf}, but with the true dynamics \eqref{eq:true_dynamics} in
place of the approximate dynamics \eqref{eq:model_dynamics}. Fix a compact
operating region $\Omega\subset\R^n$ and a time $t\in[t_0,0]$.

\begin{proposition}[Conservative superlevel-set inclusion]
\label{prop:set_inclusion}
Suppose there exists $\varepsilon_B(t)\in\R_{\ge 0}$ such that
\begin{equation}
    \sup_{\xv\in\Omega}
    \bigl|\Bt(\xv,t)-\Bm(\xv,t)\bigr|
    \le \varepsilon_B(t).
    \label{eq:certificate_error}
\end{equation}
Then
\begin{equation}
    \{\xv\in\Omega : \Bm(\xv,t)\ge \varepsilon_B(t)\}
    \subseteq
    \{\xv\in\Omega : \Bt(\xv,t)\ge 0\}.
    \label{eq:set_inclusion}
\end{equation}
\end{proposition}

\begin{proof}
Fix any $\xv\in\Omega$ such that $\Bm(\xv,t)\ge\varepsilon_B(t)$. From
\eqref{eq:certificate_error},
\[
    \Bt(\xv,t)
    \ge
    \Bm(\xv,t)-\varepsilon_B(t)
    \ge 0.
\]
Thus $\xv\in\{\xv\in\Omega:\Bt(\xv,t)\ge0\}$, which proves
\eqref{eq:set_inclusion}.
\end{proof}

Proposition~\ref{prop:set_inclusion} shows that, if a uniform bound on the
certificate error were available, then a tightened superlevel set of the
model-based CBVF would be contained in the corresponding true-dynamics
certified set. In practice, however, a deterministic bound such as
\eqref{eq:certificate_error} is difficult to estimate tightly. It must control
the full value-function error over the operating region $\Omega$, including
states and mismatch realizations that may be rare along closed-loop rollouts.

The online filter does not require such a global value-error bound. The
constraint in \eqref{eq:model_cbvf_qp} depends on the directional derivative of
the certificate along the dynamics. Therefore, instead of bounding the full
difference between $\Bt$ and $\Bm$, the proposed method calibrates the
one-sided directional mismatch that can weaken this constraint. The next
section makes this step precise by defining rollout-level mismatch scores and
using conformal prediction to convert them into a data-dependent safety buffer
for the CBVF filter.



%%%%%%%%%%%%%%%%%%%%%%%%%




\section{Conformal Calibration of the CBVF Filter}
\label{sec:conformal_calibration}

In this section, we calibrate the CBVF filter using rollout data. The goal is
to correct the part of the model mismatch that can make the safety constraint
too optimistic. We keep the model-based CBVF filter, but add a data-driven
buffer to its online constraint.

This buffer is computed using conformal prediction. Conformal prediction uses
calibration data to produce a threshold with finite-sample probabilistic
coverage \cite{angelopoulos2023conformal,shafer2008tutorial}. In our setting,
this threshold becomes the safety buffer for the CBVF filter. The calibrated
filter therefore keeps the same basic QP structure as the model-based filter,
but accounts for safety-relevant mismatch between the approximate model and the
true system.
%%%%%%%%%%%%%%%%%%%%%%%%%%%%%%%%
\subsection{One-sided mismatch score}

For online safety, the relevant quantity is not the full model mismatch, but
only the part that can make the certificate evolve less favorably than predicted
by the approximate model. We therefore use a one-sided positive score that
retains safety-deteriorating errors and ignores discrepancies for which the
true dynamics do not decrease the certificate more than the approximate model
predicts.
%%%%%%%%%%%%%5
Let $N\in\N$ denote the number of calibration rollouts, indexed by
$i\in\{1,\dots,N\}$. Let $\xv_i(s)$, $\uv_i(s)$, and $\dv_i(s)$ denote the
state, applied input, and realized disturbance along the $i$th rollout. For
each time $s\in[t_0,0]$, define
\begin{equation}
\begin{aligned}
    \eta_i(s)
    :=
    \Bigl[
    &\min_{\dv\in\Dv}
    \ip{\nabla_{\xv}\Bm(\xv_i(s),s)}
    {\fapp(\xv_i(s),\uv_i(s),\dv)} \\
    &-
    \ip{\nabla_{\xv}\Bm(\xv_i(s),s)}
    {f(\xv_i(s),\uv_i(s),\dv_i(s))}
    \Bigr]_+ .
\end{aligned}
\label{eq:eta_def}
\end{equation}

The positive-part operator in \eqref{eq:eta_def} keeps only the cases where the
approximate model predicts a more favorable certificate derivative than the
true dynamics provide. If the directional derivative along the true dynamics is
no smaller than the model prediction, then $\eta_i(s)=0$, and no additional
conservatism is introduced at time $s$.
%%%%%%%%%%%%%%
\begin{assumption}[Calibration measurements]
\label{ass:calibration_access}
For each calibration rollout $i\in\{1,\dots,N\}$, the state trajectory
$\xv_i(s)$ and the applied control $\uv_i(s)$ are available on $[t_0,0]$.
In addition, the scalar directional derivative
\[
    \ip{\nabla_{\xv}\Bm(\xv_i(s),s)}
    {f(\xv_i(s),\uv_i(s),\dv_i(s))}
\]
that appears in the score \eqref{eq:eta_def} can be computed or estimated.

In simulation or model-known settings, this quantity can be computed by
evaluating the true dynamics along the rollout, as in standard
reachability-based and CBVF computations
\cite{mitchell2005time,choi2021cbvf}. In experimental settings, it can be
estimated from logged state measurements, for example through numerical
differentiation or state estimation, or by using observer-based methods to
estimate disturbance and uncertainty effects
\cite{vanbreugel2020numerical,alan2023dobcbf,chen2023robusteso}. Equivalently,
using \eqref{eq:true_dynamics}, it is enough to estimate
\[
    \ip{\nabla_{\xv}\Bm(\xv_i(s),s)}{\dot{\xv}_i(s)}
\]
from the rollout data.
\end{assumption}

Assumption~\ref{ass:calibration_access} ensures that the pointwise score in
\eqref{eq:eta_def} can be computed from each calibration rollout. The key
requirement is not a complete model of the true dynamics, but the directional
effect of the observed trajectory on the CBVF. Once this pointwise mismatch
score is available, the next step is to convert each rollout into one scalar
calibration score.

Since each calibration example is a full trajectory, we aggregate the
pointwise score over the horizon:
\begin{equation}
    s_i:=\sup_{s\in[t_0,0]}\eta_i(s).
    \label{eq:trajectory_score}
\end{equation}
The score $s_i$ records the largest safety-degrading mismatch observed along
the $i$th rollout. Taking the supremum is conservative for an individual
trajectory, since one large value determines the rollout score. However, it
matches the finite-horizon safety objective because it controls the mismatch
over the entire interval $[t_0,0]$. It is also more targeted than the global
certificate-error bound in \eqref{eq:certificate_error}, since it calibrates
directional mismatch along sampled trajectories rather than a uniform
value-function error over the whole region $\Omega$.

The rollout-level scores in \eqref{eq:trajectory_score} are then used to
compute the split-conformal safety buffer.

%%%%%%%%%%%%%%%%%%%%%%%%%%%%%%%%
\subsection{Split-conformal buffer}

The rollout-level scores in \eqref{eq:trajectory_score} are now used to compute
a safety buffer. The goal is to choose a threshold that is large enough to cover
the safety-degrading mismatch observed in calibration, while still being
computed directly from data.

\begin{assumption}[Exchangeability]
\label{ass:exchangeability}
The calibration scores $\{s_i\}_{i=1}^N$ from \eqref{eq:trajectory_score} and
the deployment score
\begin{equation}
    s_{\mathrm{new}}
    :=
    \sup_{s\in[t_0,0]}\eta_{\mathrm{new}}(s)
    \label{eq:new_score}
\end{equation}
are exchangeable. The approximate model $\fapp$, the certificate $\Bm$, the
score construction in \eqref{eq:eta_def}--\eqref{eq:trajectory_score}, and the
nominal policy are fixed across calibration and deployment.
\end{assumption}

Assumption~\ref{ass:exchangeability} means that the calibration and deployment
scores are generated under the same closed-loop setup. In practice, this
requires collecting calibration rollouts using the same nominal policy,
approximate model, CBVF table, score definition, initial-state distribution,
and disturbance or mismatch regime used at deployment. If any of these
components changes substantially, the calibration scores no longer represent
the deployment setting and the buffer should be recalibrated.

Under Assumption~\ref{ass:exchangeability}, split conformal prediction gives a
finite-sample upper bound on the deployment score
\cite{angelopoulos2023conformal,shafer2008tutorial}. Let $\delta\in(0,1)$ be
the target miscoverage level, and let
\[
    s_{(1)}\le \cdots \le s_{(N)}
\]
denote the ordered calibration scores. For $\delta\ge 1/(N+1)$, define
\begin{equation}
    \widehat\xi_{\delta}
    :=
    s_{(k)},
    \qquad
    k:=\left\lceil (N+1)(1-\delta)\right\rceil .
    \label{eq:conformal_quantile}
\end{equation}
The condition $\delta\ge 1/(N+1)$ ensures that $k\le N$. For smaller
miscoverage levels, one may set $\widehat\xi_{\delta}=\infty$, which preserves
validity but gives a vacuous buffer.

The quantity $\widehat\xi_{\delta}\in\R_{\ge 0}$ is the split-conformal safety
buffer. Smaller values of $\delta$ require higher coverage and therefore
typically produce a larger buffer. The finite-sample guarantee holds for any
calibration size $N$ for which the quantile in
\eqref{eq:conformal_quantile} is defined. However, smaller calibration sets
give a coarser empirical quantile and can make $\widehat\xi_{\delta}$ more
variable or more conservative. Increasing $N$ usually gives a more stable
estimate of the safety buffer, but requires more calibration rollouts.

\begin{proposition}[Finite-sample conformal coverage]
\label{prop:cp_coverage}
Under Assumption~\ref{ass:exchangeability},
\begin{equation}
    \PP\!\left(s_{\mathrm{new}}\le \widehat\xi_{\delta}\right)
    \ge 1-\delta .
    \label{eq:cp_coverage}
\end{equation}
\end{proposition}

\begin{proof}
The result follows from the standard split-conformal coverage guarantee applied
to the exchangeable calibration scores in \eqref{eq:trajectory_score} and the
deployment score in \eqref{eq:new_score}
\cite{angelopoulos2023conformal,shafer2008tutorial}.
\end{proof}

Proposition~\ref{prop:cp_coverage} states that the calibrated buffer upper
bounds the deployment rollout score with probability at least $1-\delta$. This
coverage event is the key probabilistic ingredient in the safety proof. The
next section uses the calibrated buffer $\widehat\xi_{\delta}$ to define the
online CBVF filter and then shows how the coverage event leads to
finite-horizon safety for the true closed-loop system.
%%%%%%%%%%%%%%%%%%%%%%%%%%%%%%%%%%%%5

\section{Calibrated CBVF Filter and Safety Guarantee}
\label{sec:calibrated_filter_safety}

This section defines the calibrated online filter and proves its safety
guarantee. The filter keeps the same goal as the model-based CBVF filter: stay
as close as possible to the nominal policy while enforcing a safety condition.
The difference is that the safety condition is tightened using the conformal
buffer computed from rollout data.

This tightening compensates for the part of the model mismatch that can make
the approximate model too optimistic. We then show that, when the conformal
coverage event holds, the calibrated filter keeps the true trajectory inside
the safe set over the finite horizon. The full workflow, from nominal policy
construction to online calibrated filtering, is summarized at the end of the
section in Algorithm~\ref{alg:conformal_cbvf_workflow}.

\subsection{Calibrated CBVF-QP}

The conformal CBVF filter is defined by replacing the model-based inequality in
\eqref{eq:model_cbvf_qp} with a calibrated inequality:
\begin{equation}
\begin{aligned}
    \pi_{\mathrm{safe}}^{\mathrm{CP}}(\xv,t)
    =\;& \arg\min_{\uv\in\Uv}
    \norm{\uv-\pi_{\mathrm{nom}}(\xv,t)}^2 \\
    \text{s.t.}\;&
    \partial_t \Bm(\xv,t)
    + \gamma \Bm(\xv,t) \\
    &\quad+
    \min_{\dv\in\Dv}
    \ip{\nabla_{\xv}\Bm(\xv,t)}
    {\fapp(\xv,\uv,\dv)}
    \ge \widehat\xi_{\delta}.
\end{aligned}
\label{eq:cp_filter}
\end{equation}
Compared with \eqref{eq:model_cbvf_qp}, the objective is unchanged. The only
difference is the right-hand side of the CBVF constraint. The buffer
$\widehat\xi_\delta$ forces the approximate-model derivative to remain above a
calibrated threshold, so that the filter accounts for the safety-degrading
directional mismatch measured by the scores in \eqref{eq:eta_def}.

The safety proof uses two local conditions. First, the CBVF must be
differentiable along the part of the state-time space visited by the closed-loop
trajectory. Second, the calibrated constraint must admit an admissible input
whenever the trajectory is inside the certified set.

\begin{assumption}[Local regularity]
\label{ass:local_regularity}
The function $\Bm$ is continuously differentiable on the operating region
visited by the closed-loop system.
\end{assumption}

\begin{assumption}[Pointwise feasibility]
\label{ass:qp_feasibility}
For every $(\xv,t)$ visited by the closed-loop system with
$\Bm(\xv,t)\ge 0$, there exists a control input $\uv\in\Uv$ satisfying the
constraint in \eqref{eq:cp_filter}.
\end{assumption}

The feasibility condition becomes more restrictive when the buffer
$\widehat\xi_\delta$ increases. In particular, smaller values of $\delta$
typically give a larger conformal quantile in \eqref{eq:conformal_quantile},
which tightens the constraint in \eqref{eq:cp_filter}.

\begin{remark}
The CBVF in \eqref{eq:model_cbvf_vi} is a viscosity solution and need not be
smooth everywhere. Assumption~\ref{ass:local_regularity} is therefore imposed
only on the region visited by the online filter. Nonsmooth extensions could use
generalized gradients, but this paper focuses on the differentiable case.
\end{remark}

\subsection{Finite-horizon safety guarantee}

We now state the main safety result. The result combines the calibrated filter
in \eqref{eq:cp_filter} with the conformal coverage guarantee in
\eqref{eq:cp_coverage}.

\begin{theorem}[Finite-horizon high-probability safety]
\label{thm:main_safety}
Let $\Bm$ be the model-based CBVF in \eqref{eq:model_cbvf}, and let
$\xv(\cdot)$ be a closed-loop trajectory of the true system
\eqref{eq:true_dynamics} on $[t_0,0]$ under the calibrated filter
\eqref{eq:cp_filter}. Suppose Assumptions~\ref{ass:regularity},
\ref{ass:exchangeability}, \ref{ass:local_regularity}, and
\ref{ass:qp_feasibility} hold. Assume also that
\begin{equation}
    \Bm(\xv(t_0),t_0)\ge 0.
    \label{eq:initial_condition}
\end{equation}
Let $\mathcal E_{\delta}$ denote the conformal coverage event
\begin{equation}
    \mathcal E_{\delta}
    :=
    \left\{
    \sup_{s\in[t_0,0]}\eta(s)\le \widehat\xi_{\delta}
    \right\},
    \label{eq:coverage_event}
\end{equation}
where $\eta(s)$ is the one-sided mismatch score in \eqref{eq:eta_def}
evaluated along the deployment trajectory. Then, on
$\mathcal E_{\delta}$,
\begin{equation}
    \Bm(\xv(s),s)\ge 0,
    \qquad \forall s\in[t_0,0].
    \label{eq:cbvf_nonnegative}
\end{equation}
Therefore,
\begin{equation}
    \xv(s)\in\safe,
    \qquad \forall s\in[t_0,0].
    \label{eq:true_safety}
\end{equation}
Moreover,
\begin{equation}
    \PP\!\left(
    \xv(s)\in\safe \text{ for all } s\in[t_0,0]
    \right)
    \ge 1-\delta .
    \label{eq:prob_safe}
\end{equation}
\end{theorem}

\begin{proof}
By the definition of the one-sided mismatch score in \eqref{eq:eta_def}, for
almost every $s\in[t_0,0]$ along the deployment trajectory,
\begin{equation}
\begin{aligned}
    &\ip{\nabla_{\xv}\Bm(\xv(s),s)}
    {f(\xv(s),\uv(s),\dv(s))}  \\
    &\ge
    \min_{\dv\in\Dv}
    \ip{\nabla_{\xv}\Bm(\xv(s),s)}
    {\fapp(\xv(s),\uv(s),\dv)}
    -\eta(s).
\end{aligned}
\label{eq:eta_ineq}
\end{equation}
Since the control input is generated by the calibrated filter
\eqref{eq:cp_filter},
\begin{equation}
\begin{aligned}
    &\partial_t \Bm(\xv(s),s)
    + \gamma \Bm(\xv(s),s) \\
    &\quad+
    \min_{\dv\in\Dv}
    \ip{\nabla_{\xv}\Bm(\xv(s),s)}
    {\fapp(\xv(s),\uv(s),\dv)}
    \ge \widehat\xi_{\delta}.
\end{aligned}
\label{eq:filter_ineq}
\end{equation}
Combining \eqref{eq:eta_ineq} and \eqref{eq:filter_ineq}, and using
Assumption~\ref{ass:local_regularity}, gives
\begin{equation}
    \frac{d}{ds}\Bm(\xv(s),s)
    +\gamma \Bm(\xv(s),s)
    \ge \widehat\xi_{\delta}-\eta(s)
    \quad \text{a.e. on } [t_0,0].
    \label{eq:diff_ineq_pre}
\end{equation}
On the event $\mathcal E_{\delta}$ in \eqref{eq:coverage_event},
$\eta(s)\le \widehat\xi_{\delta}$ for all $s\in[t_0,0]$. Hence
\begin{equation}
    \frac{d}{ds}\Bm(\xv(s),s)
    +\gamma \Bm(\xv(s),s)
    \ge 0
    \quad \text{a.e. on } [t_0,0].
    \label{eq:diff_ineq}
\end{equation}
Since $\Bm$ is continuously differentiable on the visited region and the true
trajectory is absolutely continuous by Assumption~\ref{ass:regularity}, the map
$s\mapsto \Bm(\xv(s),s)$ is absolutely continuous. Multiplying
\eqref{eq:diff_ineq} by the integrating factor $e^{\gamma s}$ and integrating
from $t_0$ to $s$ yields
\begin{equation}
    \Bm(\xv(s),s)
    \ge
    e^{-\gamma(s-t_0)}\Bm(\xv(t_0),t_0),
    \qquad \forall s\in[t_0,0].
    \label{eq:cbvf_lb}
\end{equation}
The right-hand side is nonnegative by \eqref{eq:initial_condition}, which proves
\eqref{eq:cbvf_nonnegative}. Since $\Bm(\xv,s)\le l(\xv)$ by
\eqref{eq:Bm_le_l}, we have $l(\xv(s))\ge 0$ for all $s\in[t_0,0]$. Thus
$\xv(s)\in\safe$ for all $s\in[t_0,0]$, proving \eqref{eq:true_safety}.

Finally, Proposition~\ref{prop:cp_coverage} and
Assumption~\ref{ass:exchangeability} imply that $\mathcal E_{\delta}$ holds
with probability at least $1-\delta$. Since safety holds on
$\mathcal E_{\delta}$, \eqref{eq:prob_safe} follows.
\end{proof}

\begin{remark}
Theorem~\ref{thm:main_safety} gives a finite-horizon, trajectory-level, marginal
probabilistic safety guarantee. The probability in \eqref{eq:prob_safe} is over
the calibration rollouts and the deployment rollout. It is not a worst-case
guarantee over all possible deployment conditions.
\end{remark}

\begin{remark}[Sampled-data implementation]
\label{rem:sampled_data}
The proof is stated in continuous time, while the implemented filter is applied
at the decision times in \eqref{eq:decision_times} with zero-order hold. This
sampled-data gap can be reduced by using a faster update rate or by adding an
inter-sample safety margin
\cite{fisac2019general,alshiekh2018safe,lu2025adaptiveshielding}.
\end{remark}

%%%%%%%%%%%%%%%%%%%%%%%%%%%%%%
\begin{algorithm}[t]
\caption{Conformal CBVF shielding workflow}
\label{alg:conformal_cbvf_workflow}
\scriptsize
\begin{algorithmic}[1]
\State Train the nominal policy $\pi_{\mathrm{nom}}$ in
\eqref{eq:nominal_policy}.
\State Compute the model-based CBVF $\Bm$ from \eqref{eq:model_cbvf}.
\State Fix the approximate model $\fapp$, the CBVF $\Bm$, and
$\pi_{\mathrm{nom}}$.
\State Collect calibration rollouts $\{\xv_i(\cdot),\uv_i(\cdot)\}_{i=1}^N$.
\State Compute the pointwise scores $\eta_i(s)$ using \eqref{eq:eta_def}.
\State Form rollout scores $s_i=\sup_{s\in[t_0,0]}\eta_i(s)$ as in
\eqref{eq:trajectory_score}.
\State Compute the conformal buffer $\widehat\xi_\delta$ using
\eqref{eq:conformal_quantile}.
\State At each decision time $t_j$, solve the calibrated QP
\eqref{eq:cp_filter}.
\State Apply the filtered input on $[t_j,t_{j+1})$ with zero-order hold.
\State Safety follows from Theorem~\ref{thm:main_safety}.
\end{algorithmic}
\end{algorithm}

%%%%%%%%%%%%%%%%%%%%%%%%%%%%%%%%%%%
\section{Implementation and Experimental Validation}
\label{sec:experiments}

This section evaluates the calibrated CBVF filter on a planar reach-avoid task
with controlled model mismatch. The goal is to test whether the conformal
buffer improves safety when the online filter is built from an approximate
dynamics model.

The system is a Dubins car with state $\xv=(p_x,p_y,\theta)$ and steering
input $\uv\in[-1,1]$. The true dynamics are
\begin{equation}
    \dot p_x = v\cos\theta,\qquad
    \dot p_y = v\sin\theta,\qquad
    \dot\theta = \beta u,
    \label{eq:dubins_dynamics_exp}
\end{equation}
with forward speed $v=0.60$. The robot starts near
$(-2,-2,\pi/4)$, with uniform position jitter of $\pm 0.30$ and heading jitter
of $\pm 0.15$ rad. The goal region is centered at $(2,2)$ with radius $0.30$.
The unsafe set is a circular obstacle centered at the origin, with combined
robot-obstacle radius $r=0.70$. Safety is encoded by the signed distance to the
obstacle boundary,
\begin{equation}
    l(\xv)=\|p-p_{\mathrm{obs}}\|_2-r .
    \label{eq:signed_distance_exp}
\end{equation}

The nominal controller is a SAC-Lagrangian policy trained in the aligned
no-mismatch environment. The offline CBVF is recomputed for the reported
experiment on an $81\times 81\times 61$ spatial grid, with lookback horizon
$T=20\,\mathrm{s}$, time step $\Delta t=0.10\,\mathrm{s}$, and discount
$\gamma=0.1$. The conformal miscoverage level is $\delta=0.05$, giving a
$95\%$ split-conformal coverage level, and the calibration set contains
$N=500$ rollouts. Results are reported as mean $\pm$ standard deviation over
five random seeds, with $100$ evaluation episodes per seed.

The experiments use a matched-training protocol: a nominal policy trained under
a given dynamics setting is evaluated only under the same deployment setting.
This section reports the aligned no-mismatch case, where the model used by the
filter matches the true dynamics,
$\beta_{\mathrm{model}}=\beta_{\mathrm{truth}}=0.50$. The same protocol can be
applied separately to hard and easy mismatch policies once those policies are
trained.

We compare three methods. The first is the unshielded nominal RL policy. The
second is a least-restrictive CBVF baseline, denoted CBVF (LR), which activates
only when $\Bm(\xv,t)\le 0$ and uses a zero right-hand side in the QP
constraint. The third is the calibrated CP-CBVF filter in \eqref{eq:cp_filter},
which activates when $\Bm(\xv,t)\le 0.10$ and uses the conformal buffer
$\widehat\xi_\delta$.

\begin{table}[t]
\centering
\caption{Reach-avoid performance in the aligned no-mismatch setting. Values are mean $\pm$ standard deviation over five seeds. CBVF activates at $\Bm\le 0$; CP-CBVF activates at $\Bm\le 0.10$.}
\label{tab:experiment_results}
\scriptsize
\setlength{\tabcolsep}{2.5pt}
\begin{tabular}{llccccc}
\hline
Scenario & Method & Unsafe & Goal & Return & Interv. & $\widehat\xi_\delta$ \\
\hline
Aligned & Nominal
& $0.252 \pm 0.051$
& $0.748 \pm 0.051$
& $-520 \pm 10$
& $0$
& -- \\
Aligned & CBVF (LR)
& $0.252 \pm 0.051$
& $0.748 \pm 0.051$
& $-520 \pm 10$
& $0.3 \pm 0.1$
& $0.000$ \\
Aligned & \textbf{CP-CBVF}
& $\mathbf{0.000 \pm 0.000}$
& $\mathbf{1.000 \pm 0.000}$
& $-570 \pm 2$
& $8.4 \pm 0.7$
& $0.000$ \\
\hline
\end{tabular}
\end{table}

Table~\ref{tab:experiment_results} summarizes the aligned no-mismatch results.
The nominal policy reaches the goal in most episodes, but it still produces
unsafe rollouts. The least-restrictive CBVF baseline behaves almost identically
to the nominal policy because it activates only at the certificate boundary and
makes very few corrections. In contrast, CP-CBVF achieves zero unsafe episodes
and a $100\%$ goal-reaching rate in the aligned setting.

\begin{figure}[t]
    \centering
    \includegraphics[width=\columnwidth]{safety_comparison.png}
    \caption{Safety and task-completion comparison in the aligned no-mismatch setting. Left: unsafe episode rate. Right: goal-reaching rate. The nominal policy and least-restrictive CBVF baseline are nearly identical, while CP-CBVF removes all violations and reaches the goal in all evaluation episodes. Error bars show $\pm$ one standard deviation over five seeds.}
    \label{fig:safety_comparison}
\end{figure}

Figure~\ref{fig:safety_comparison} visualizes the same safety and goal-reaching
metrics. The nominal and CBVF (LR) bars overlap because the least-restrictive
filter rarely changes the nominal action. CP-CBVF removes the observed
violations while preserving goal completion, showing that earlier activation is
effective even when the approximate and true dynamics are aligned.

\begin{figure}[t]
    \centering
    \includegraphics[width=0.85\columnwidth]{trajectories_AlignedModel.png}
    \caption{Representative aligned-setting rollout trajectories from seed $7$.
    The black circle is the obstacle, the green circle is the goal region, and
    the star marks the nominal start. The legend distinguishes the nominal,
    CBVF (LR), and CP-CBVF trajectories. The nominal and CBVF trajectories may
    overlap on safe episodes because the least-restrictive CBVF baseline rarely
    changes the nominal action.}
    \label{fig:representative_trajectories}
\end{figure}

The representative trajectories in Fig.~\ref{fig:representative_trajectories}
show how the methods behave in the state space. The nominal policy can take a
short path toward the goal that intersects the obstacle. The CBVF (LR) baseline
follows nearly the same route because it intervenes only near the certificate
boundary. CP-CBVF instead steers around the obstacle and reaches the goal in the
shown aligned-setting rollouts.

In the aligned setting, the calibration scores are zero, so
$\widehat\xi_\delta=0$. Therefore, the improvement does not come from a nonzero
model-mismatch buffer. It comes from activating the CP-CBVF filter earlier,
before the certificate reaches zero. This behavior is consistent with the
matched no-mismatch setting, where the approximate model is not optimistic
relative to the true dynamics.

\begin{figure}[t]
    \centering
    \includegraphics[width=0.95\columnwidth]{interventions.png}
    \caption{Mean shield interventions per rollout in the aligned no-mismatch setting. CBVF (LR) makes near-zero interventions because it only activates at the certificate boundary. CP-CBVF intervenes earlier to maintain the activation margin and enforce the calibrated safety constraint.}
    \label{fig:interventions}
\end{figure}

Figure~\ref{fig:interventions} reports the number of filter interventions per
rollout. CP-CBVF intervenes more often than CBVF (LR), which is expected
because CP-CBVF activates before the certificate reaches zero. The added
interventions produce a performance tradeoff: the mean CP-CBVF return is lower
than the nominal return, reflecting the longer but safe route around the
obstacle.

The empirical results are consistent with the conformal safety guarantee.
Proposition~\ref{prop:cp_coverage} gives
\[
    \mathbb{P}\!\left(s_{\mathrm{new}}\le \widehat\xi_\delta\right)
    \ge 1-\delta = 0.95,
\]
and Theorem~\ref{thm:main_safety} converts this coverage event into a
finite-horizon safety statement for the calibrated filter. Across $500$
deployment rollouts in the aligned setting, CP-CBVF has empirical unsafe rate
$0.000$. Overall, the aligned experiment shows that CP-CBVF improves safety
while preserving task completion under the matched no-mismatch deployment
protocol.

%%%%%%%%%%%%%%%%%%%%%%%%%%%%%%%%

\section{Conclusion}
% To be added.

\bibliographystyle{IEEEtran}
\bibliography{conformal_cbvf_shielding}

\end{document}